\chapter{Introduction}
\label{ch:intro}

\section{The question}

A ride-on mower converted into a robot has to mow a whole lawn on its own.
Which path should it follow? A good path cuts every square metre of grass,
wastes little time driving over grass already cut, never leaves the lawn or
hits a flower bed, and is one the machine can actually drive.

The last condition is where the difficulty lies. A ride-on mower steers like
a car: it cannot turn on the spot, its tightest turn has a radius of more than
a metre, and to go back along the next strip it must either make a wide loop
or a manoeuvre with a reverse leg. Most simple models of coverage ignore this:
they cut the lawn into square cells and let the mower jump from a cell to its
neighbour, turning by 90 degrees whenever it likes. Such a model is a fine
proof of concept for comparing algorithms, but the paths it produces cannot be
driven by a car-like machine.

This course builds the realistic version step by step.

\section{From a grid to a real lawn}

\begin{center}
\begin{tabular}{@{}lll@{}}
  \toprule
  & grid model & this course \\
  \midrule
  lawn & square cells & polygon in metres, with holes \\
  mower & a cell and one of 8 headings & pose $(x, y, \theta)$, real
    dimensions \\
  motion & jump to a neighbouring cell & arcs of curvature
    $|\kappa| \le \kmax$, forward or reverse \\
  mowed area & the cells visited & the area swept by the blades \\
  cost & number of moves and turns & time, missed area, area cut twice \\
  lower bound & $N - 1$ moves & area divided by an effective width \\
  \bottomrule
\end{tabular}
\end{center}

\section{What this course covers}

This first edition lays the foundations: what a lawn is, how the mower moves,
how to join two poses by the shortest drivable path, and how to measure a
path. Each chapter derives its results completely, then shows where they are
in the code.

\begin{description}
  \item[Chapter~\ref{ch:lawn}, the lawn.]
    \leavevmode\\
    Polygons with holes, their area, the point-in-polygon test, and the
    raster on which coverage is measured, with a bound on its error.
  \item[Chapter~\ref{ch:mower}, the mower.]
    \leavevmode\\
    The machine seen from above, the kinematic bicycle model derived from
    rolling without slipping, Ackermann steering, the exact motion along an
    arc, and the area the blades sweep in a turn.
  \item[Chapter~\ref{ch:paths}, shortest paths.]
    \leavevmode\\
    Dubins paths (forward only) and Reeds--Shepp paths (with reverse), and the
    case every mowing pattern meets: turning into the next strip.
  \item[Chapter~\ref{ch:measuring}, measuring a path.]
    \leavevmode\\
    Coverage, collisions, time, and a lower bound on the length of any path
    that mows the whole lawn. A first experiment shows why the border of the
    lawn needs special treatment.
\end{description}

The following editions will build planners on these foundations: headland
passes around the border, parallel swaths joined by feasible turns, the
decomposition of complex lawns into simpler regions and the order in which to
mow them, and finally learning and search to make the high-level choices.

\section{The code}

Everything in this course is implemented in the Python package
\code{mowersim}, in the same repository. The figures that show trajectories
and coverage maps are drawn by the code (\code{tools/course\_figures.py}), so
they show what the code computes. A single function, \code{evaluate}, measures
every trajectory, whatever produced it: two methods can then only differ by
their paths, never by the way they are scored.

\section{Conventions}

\begin{itemize}
  \item Lengths are in metres, times in seconds, angles in radians in the
    formulas (degrees only when quoting a value).
  \item The \emph{pose} of the mower is the position of the middle of its rear
    axle and its heading (Chapter~\ref{ch:mower} explains this choice).
  \item Angles are counter-clockwise, and a positive curvature turns to the
    left.
  \item The numerical examples use the default mower of the code: a
    tractor-style ride-on mower with a cutting width of
    \SI{0.97}{\metre} and a minimum turning radius of \SI{1.43}{\metre}. All
    its dimensions are parameters (Appendix~\ref{ch:appendix}).
\end{itemize}
